\hmsubsection{Object Detection}{OAH}{}

This section describes the object detection stage of the vision system,
which is responsible for identifying the colour, position, and size of
balls present in each processed frame. Detection is based on a combination
of colour segmentation in the HSV colour space, Hough Circle Transform,
geometric contour validation, ensemble merging of detections from multiple
methods, and persistent object tracking across frames. The following
subsections detail each component in turn.

\hmsubsubsection{HSV Colour Segmentation}{OAH}{}

Colour-based object detection is performed using thresholding in the
Hue-Saturation-Value (HSV) colour space. HSV is preferred over BGR for this
application because it separates chromatic information (hue) from illumination
(value), making the thresholds less sensitive to variations in lighting conditions.

For each frame, separate binary masks are produced for red and blue objects. Each
colour is defined by multiple HSV ranges, which are applied independently and
combined using a bitwise OR operation. The resulting mask covers all pixels whose
HSV values fall within any of the defined ranges.

Red requires two pairs of ranges due to the circular nature of the hue axis in
OpenCV, where hue values are represented in the interval $[0, 179]$. Red occupies
both ends of this interval, wrapping around at the boundary. The first pair targets
normally illuminated red surfaces: $H \in [160, 179]$ and $H \in [0, 10]$, both
with $S \in [120, 255]$ and $V \in [40, 255]$. The second pair targets dark red
surfaces in shadow or under low ambient light: $H \in [160, 179]$ and
$H \in [0, 10]$, both with $S \in [80, 255]$ and $V \in [20, 180]$. The lower
saturation and value floors in the shadow pair allow detection of the same physical
ball when it is partially occluded or shaded by the robot arm.

Blue is defined by three ranges. The first range, $H \in [85, 115]$ with
$S \in [55, 255]$ and $V \in [60, 255]$, captures light and cyan-tinted blue
surfaces. The second range, $H \in [100, 135]$ with $S \in [90, 255]$ and
$V \in [40, 255]$, captures standard saturated blue surfaces. A third range,
$H \in [95, 135]$ with $S \in [70, 255]$ and $V \in [20, 190]$, captures dark
navy-blue surfaces where saturation and brightness are significantly lower.

Following mask construction, morphological operations are applied to reduce noise
and close gaps in the mask. An opening operation with a $3 \times 3$ elliptical
kernel removes isolated noise pixels, and a closing operation with a
$13 \times 13$ elliptical kernel connects fragmented regions caused by specular
highlights on the ball surface. Finally, a flood fill from the image corner is
used to fill any remaining interior holes in the mask before contour extraction.

\begin{figure}[h]
\centering
\begin{tikzpicture}[scale=0.85]

% --- Hue circle ---
\foreach \angle in {0,2,...,358} {
    \pgfmathsetmacro{\hsvhue}{\angle/360}
    \definecolor{huecolor}{hsb}{\hsvhue, 1, 1}
    \fill[huecolor] (0,0) -- (\angle:1.5) arc (\angle:\angle+2:1.5) -- cycle;
}
\fill[white] (0,0) circle (0.8);
\node[font=\small\bfseries] at (0,0) {Hue};

% Red range 1 normal: H=160-179
\fill[red!60, opacity=0.5] (0,0) -- ({160*2}:1.5) arc ({160*2}:{179*2}:1.5) -- cycle;
% Red range 2 normal: H=0-10
\fill[red!60, opacity=0.5] (0,0) -- (0:1.5) arc (0:{10*2}:1.5) -- cycle;
% Blue range union: H=85-135
\fill[blue!60, opacity=0.5] (0,0) -- ({85*2}:1.5) arc ({85*2}:{135*2}:1.5) -- cycle;

% Labels
\node[font=\tiny, red] at ({170*2}:1.90) {$[160,179]$};
\node[font=\tiny, red] at ({5*2}:2.10) {$[0,10]$};
\node[font=\tiny, blue] at ({110*2}:1.95) {$[85,135]$};
\node[font=\tiny] at (0, -2.1) {$H = 0/180$};
\draw[dashed, gray] (0,0) -- (0:1.7);

\node[font=\small\bfseries] at (0, 2.5) {HSV Hue axis ($0$--$179$)};

% --- Arrow ---
\draw[->, thick] (2.3, 0) -- (3.3, 0);

% --- Mask diagram ---
% BGR frame box
\draw[thick] (3.5,-1.2) rectangle (6.5,1.2);
\fill[gray!15] (3.5,-1.2) rectangle (6.5,1.2);
\fill[red!70]  (4.4,-0.3) circle (0.55);
\fill[blue!70] (5.6, 0.2) circle (0.45);
\node[font=\tiny] at (5.0,-1.5) {BGR frame};

% Arrow
\draw[->, thick] (6.7, 0) -- (7.7, 0);
\node[font=\tiny, above] at (7.2, 0.05) {mask};

% Binary mask box
\draw[thick] (7.9,-1.2) rectangle (10.9,1.2);
\fill[black!85] (7.9,-1.2) rectangle (10.9,1.2);
\fill[white]  (8.8,-0.3) circle (0.55);
\fill[white]  (10.0, 0.2) circle (0.45);
\node[font=\tiny, white] at (9.4,-1.5) {Binary mask};

\end{tikzpicture}
\caption[HSV hue axis and resulting binary masks for red and blue detection.]{Left: the HSV hue axis ($0$--$179$) shown as a colour wheel. Red occupies
both ends of the axis and is therefore represented by two pairs of ranges: one pair
for normally illuminated surfaces ($S \geq 120$) and one pair for surfaces in shadow
($S \geq 80$). Both pairs share the same hue intervals $[160, 179]$ and $[0, 10]$.
Blue occupies a continuous region and is represented by three ranges covering
$H \in [85, 135]$: one for light and cyan-tinted blue surfaces, one for standard
saturated blue, and one for dark navy-blue surfaces.
Right: the resulting binary masks isolate only the relevant colour regions from the
BGR frame.}
\label{fig:hsv_mask}
\end{figure}

\hmsubsubsection{Hough Circle Transform}{OAH}{}

In addition to HSV colour segmentation, the system employs the Hough Circle
Transform as a complementary geometric detection method. While HSV segmentation
detects objects based on colour, the Hough Circle Transform detects circular shapes
directly from edge gradients in the grayscale image, independently of colour
information. This makes the two methods sensitive to different properties of the
same object, which is the foundation of the ensemble approach described in
Section~8.3.4.

The transform is applied to the Gaussian-blurred grayscale frame using
OpenCV's \texttt{HoughCircles} function with the gradient method. The accumulator
resolution is set to $dp = 1$, meaning the accumulator has the same resolution as
the input image. The minimum distance between detected circle centres is set to
$3\times$ the minimum radius, or 20 pixels, whichever is larger. This prevents
multiple detections of the same ball while permitting two balls to be detected side
by side. The edge detection threshold is set to 50 and the accumulator threshold to
35, which is permissive enough to detect matte ball surfaces while relying on the
subsequent colour and shape validation steps to reject false positives.

Each candidate circle returned by the transform is subjected to two additional
validation steps before being accepted. First, the colour of the detected region is
determined by sampling HSV values within 60\% of the circle radius from its centre,
excluding specular highlight pixels where $S < 30$ and $V > 210$. Candidates whose
colour cannot be classified as red or blue are rejected. Second, a local shape
verification is performed within $1.5\times$ the circle radius. A local HSV mask is
constructed and its largest contour is approximated to a polygon. Candidates whose
contour approximates to six or fewer vertices are rejected as non-circular objects
such as cubes or rectangular surfaces. An additional Intersection over Union (IoU)
check compares the colour mask against an ideal circle of the same radius; candidates
with $\mathrm{IoU} < 0.45$ are rejected as irregular objects. A fixed shape
confidence of 0.88 is assigned to all accepted Hough detections, reflecting that
geometric circle fitting provides strong but not perfect shape verification.
The final detection confidence is computed as a weighted combination of the
shape and colour scores:
\begin{equation}
    c_{\mathrm{raw}} = 0.80 \cdot c_{\mathrm{shape}} + 0.20 \cdot c_{\mathrm{colour}}
\end{equation}
and the final confidence is taken as the square root to compress the distribution
towards~1:
\begin{equation}
    c_{\mathrm{conf}} = \sqrt{\,c_{\mathrm{raw}}}
\end{equation}

\begin{figure}[h]
\centering
\begin{tikzpicture}[scale=0.8]

% --- Step 1: Grayscale image with circle edge ---
\draw[thick] (0,0) rectangle (3,3);
\fill[gray!15] (0,0) rectangle (3,3);
\fill[gray!35] (1.5,1.5) circle (1.1);
\draw[very thick, black] (1.5,1.5) circle (1.1);
\node[font=\small\bfseries] at (1.5, 3.5) {Grayscale frame};
\node[font=\footnotesize] at (1.5, -0.5) {1.\ Edge detection};

% --- Arrow ---
\draw[->, thick] (3.3, 1.5) -- (4.2, 1.5);

% --- Step 2: Accumulator heatmap ---
\draw[thick] (4.4, 0) rectangle (7.4, 3);
\fill[gray!15] (4.4,0) rectangle (7.4,3);

% 5x5 heatmap grid, centre at (5.9, 1.5), cell size 0.5
\foreach \i/\j/\intensity in {
    0/0/10, 1/0/20, 2/0/15, 3/0/10, 4/0/5,
    0/1/20, 1/1/40, 2/1/50, 3/1/35, 4/1/15,
    0/2/15, 1/2/55, 2/2/95, 3/2/50, 4/2/20,
    0/3/20, 1/3/40, 2/3/45, 3/3/30, 4/3/10,
    0/4/10, 1/4/15, 2/4/20, 3/4/15, 4/4/5
}{
    \fill[blue!\intensity] ({4.4+\i*0.6}, {\j*0.6}) rectangle ({4.4+\i*0.6+0.6}, {\j*0.6+0.6});
    \draw[gray!40] ({4.4+\i*0.6}, {\j*0.6}) rectangle ({4.4+\i*0.6+0.6}, {\j*0.6+0.6});
}
% Peak cell
\fill[red!85] ({4.4+2*0.6}, {2*0.6}) rectangle ({4.4+3*0.6}, {3*0.6});
\node[font=\tiny, white, align=center] at ({4.4+2.5*0.6}, {2.5*0.6}) {peak};

\node[font=\small\bfseries] at (5.9, 3.5) {Accumulator};
\node[font=\footnotesize] at (5.9, -0.5) {2.\ Vote accumulation};

% --- Arrow ---
\draw[->, thick] (7.6, 1.5) -- (8.5, 1.5);

% --- Step 3: Detected circle ---
\draw[thick] (8.7,0) rectangle (11.7,3);
\fill[gray!15] (8.7,0) rectangle (11.7,3);
\fill[gray!35] (10.2,1.5) circle (1.1);
\draw[very thick, red!80, dashed] (10.2,1.5) circle (1.1);
\fill[red!80] (10.2,1.5) circle (0.08);
\draw[red!70, thick] (10.2,1.5) -- (11.3,1.5);
\node[font=\small] at (10.8, 1.7) {$r$};
\node[font=\small\bfseries] at (10.2, 3.5) {Detected circle};
\node[font=\footnotesize] at (10.2, -0.5) {3.\ Peak = circle};

\end{tikzpicture}
\caption[Hough Circle Transform in three steps.]{Hough Circle Transform in three
steps. Edge pixels cast votes in an accumulator matrix for all circles that could
have produced them. A peak in the accumulator (red cell) corresponds to a circle
centre and radius supported by many edge pixels.}
\label{fig:hough}
\end{figure}


\hmsubsubsection{Geometric Validation}{OAH}{}

Following HSV mask construction, the system extracts external contours from the
binary mask using OpenCV's contour detection. Each contour is subjected to a
series of geometric gate filters to determine whether it corresponds to a spherical
object. A contour that fails any of the following criteria is rejected.

The minimum enclosed circle of the contour is computed to obtain the radius, which
must fall within the configured range of 10 to 150 pixels at full scale, scaled
proportionally to the processing resolution. The centroid is computed from the
first-order image moments rather than the enclosing circle centre, as moment-based
centroids are more stable under partial occlusion and surface reflections.

Circularity is computed as:
\begin{equation}
    C = \frac{4\pi A}{P^2}
\end{equation}
where $A$ is the contour area and $P$ is the perimeter. A perfect circle yields
$C = 1.0$, while a perfect square yields $C = \pi/4 \approx 0.785$. The minimum
accepted circularity is $C \geq 0.65$ for contours that lie fully within the frame.
For contours touching the image boundary, which may represent physically clipped
balls producing a D-shape, the threshold is relaxed to $C \geq 0.55$.

Following the circularity check, each contour is approximated to a polygon using a
Douglas--Peucker approximation with an epsilon equal to 2\% of its perimeter. A
circular cross-section produces eight or more vertices at this tolerance, whereas
rectangles and cubes produce four to six. Contours that approximate to six or fewer
vertices are rejected for normal contours, and four or fewer for edge-clipped
contours.

The aspect ratio of the bounding rectangle must satisfy
$\min(w,h)/\max(w,h) \geq 0.75$ for normal contours and $\geq 0.65$ for
edge-clipped contours, ensuring the detected region is approximately square as
expected for a circular cross-section.

Solidity, defined as the ratio of the contour area to its convex hull area, must
be at least 0.75 to reject concave or irregular shapes.

Colour saturation is sampled from the HSV frame using a filled contour mask,
excluding specular highlight pixels where $S < 30$ and $V > 210$. For red
contours, a mean saturation score below 0.40 results in unconditional rejection.
Blue contours are rejected if the saturation score falls below 0.28, or below
0.40 when fewer than 65\% of the non-specular pixels within the contour lie
within the blue hue range $H \in [85, 135]$. This asymmetry reflects that
light and cyan-tinted blue surfaces can exhibit lower mean saturation than red
surfaces while still being valid detections, provided the hue distribution
confirms a blue object.

A confidence score is assigned to each accepted detection based on how far each
geometric and colour metric exceeds its base threshold:
\begin{equation}
    \mathrm{confidence} = 0.70 + 0.30 \cdot (0.40 \cdot C_{\mathrm{bonus}}
    + 0.20 \cdot A_{\mathrm{bonus}} + 0.20 \cdot S_{\mathrm{bonus}}
    + 0.20 \cdot K_{\mathrm{bonus}})
\end{equation}
where
$C_{\mathrm{bonus}} = \mathrm{clip}\!\left(\frac{C - 0.65}{0.35}, 0, 1\right)$,
$A_{\mathrm{bonus}} = \mathrm{clip}\!\left(\frac{A_r - 0.75}{0.25}, 0, 1\right)$,
$S_{\mathrm{bonus}} = \mathrm{clip}\!\left(\frac{\sigma - 0.75}{0.25}, 0, 1\right)$,
and
$K_{\mathrm{bonus}} = \mathrm{clip}\!\left(\frac{k - 0.30}{0.70}, 0, 1\right)$
are the normalised bonuses for circularity, aspect ratio, solidity, and colour
saturation, respectively, each in the range $[0, 1]$. This formulation yields a
minimum confidence of 0.70 for any detection that passes all gate filters, and a
maximum of 1.00 for a geometrically ideal, fully saturated ball.

\begin{figure}[h]
\centering
\begin{tikzpicture}[
    box/.style={draw, thick, rounded corners=3pt, minimum width=2.3cm,
                minimum height=0.9cm, align=center, font=\small},
    arr/.style={->, thick}
]

% --- Row of shapes to test ---
% Ball (passes)
\fill[red!50] (0,0) circle (0.6);
\node[font=\footnotesize] at (0,-1.0) {Ball};
\node[font=\footnotesize, green!60!black] at (0,-1.4) {$\surd$ pass};

% Cube (fails circularity)
\fill[gray!40] (2.5,-0.5) rectangle (3.5,0.5);
\node[font=\footnotesize] at (3.0,-1.0) {Cube};
\node[font=\footnotesize, red!70] at (3.0,-1.4) {\texttimes\ fail};

% Elongated (fails aspect ratio)
\fill[gray!40] (5.0,-0.2) rectangle (6.6,0.2);
\node[font=\footnotesize] at (5.8,-1.0) {Elongated};
\node[font=\footnotesize, red!70] at (5.8,-1.4) {\texttimes\ fail};

% Concave (fails solidity)
\fill[gray!40] (7.8,0.5) -- (8.8,0.5) -- (8.3,-0.1) -- (9.3,0.5) -- (9.3,-0.5) -- (7.8,-0.5) -- cycle;
\node[font=\footnotesize] at (8.55,-1.0) {Concave};
\node[font=\footnotesize, red!70] at (8.55,-1.4) {\texttimes\ fail};

% --- Gate labels above ---
\draw[dashed, gray] (1.5, 1.2) -- (1.5,-1.7);
\draw[dashed, gray] (4.3, 1.2) -- (4.3,-1.7);
\draw[dashed, gray] (7.0, 1.2) -- (7.0,-1.7);

\node[font=\footnotesize, align=center] at (0,   1.0) {$C \geq 0.65$};
\node[font=\footnotesize, align=center] at (3.0, 1.0) {vertices $> 6$};
\node[font=\footnotesize, align=center] at (5.8, 1.0) {aspect $\geq 0.75$};
\node[font=\footnotesize, align=center] at (8.55,1.0) {solidity $\geq 0.75$};

\end{tikzpicture}
\caption[Geometric gate filters applied to each contour.]{Geometric gate filters
applied to each contour. A spherical ball (left) passes all four criteria:
circularity $C \geq 0.65$, more than six polygon vertices, aspect ratio $\geq 0.75$,
and solidity $\geq 0.75$. Non-spherical objects such as cubes, elongated shapes,
and concave regions are rejected at the corresponding gate. For balls touching the
image boundary, all thresholds are relaxed to tolerate the D-shape produced by
physical clipping.}
\label{fig:geometric_gates}
\end{figure}

\hmsubsubsection{Ensemble Detection and Confidence Scoring}{OAH}{}

The system combines detections from multiple sources using an ensemble approach
before producing the final output. All detections from HSV segmentation and Hough
Circle Transform are collected and passed through a two-stage merging and filtering pipeline: a Union-Find cluster merge followed by a post-merge Non-Maximum Suppression (NMS) pass to eliminate duplicate detections.

In the first stage, detections are grouped into clusters using a Union-Find
structure. Two detections are placed in the same cluster if they share the same
colour label and their centres are within $1.5\times$ the larger of their two
radii. This threshold accounts for the positional difference between HSV
moment-based centroids and Hough geometric centres, which can differ by 10--20
pixels on a ball of radius 40 pixels. Detections of different colours are never
merged.

For each cluster, the detection with the highest confidence score is selected as
the representative, and its centre is used directly as the cluster centre. The
radius is computed as the mean radius across all detections in the cluster. If the
cluster contains detections from more than one method, the confidence is boosted by
8\% per additional method, capped at 1.0:
\begin{equation}
    c_{\mathrm{final}} = \min\!\left(c_{\mathrm{best}} \cdot
    \left(1 + 0.08 \cdot (n_{\mathrm{methods}} - 1)\right),\; 1.0\right)
\end{equation}
where $c_{\mathrm{best}}$ is the highest confidence in the cluster and
$n_{\mathrm{methods}}$ is the number of distinct detection methods represented.
The detection method is labelled as ensemble when more than one method contributed.

In the second stage, a Non-Maximum Suppression (NMS) pass is applied per colour.
Detections are sorted by confidence in descending order, and any candidate whose
centre lies within $1.3\times$ the larger of the two radii of an already accepted
detection of the same colour is discarded as a duplicate. Finally, the output is
limited to a configurable maximum number of detections per colour.

During development, a Support Vector Machine (SVM) classifier was evaluated as an
additional colour verification step. The SVM was trained on HSV histograms of
cropped ball regions and applied to correct potentially mislabelled detections.
However, the classifier was found to erroneously relabel dark navy-blue balls as
red with high confidence, degrading overall accuracy. The SVM step was therefore
disabled, and the HSV and Hough ensemble alone was found to be sufficiently
reliable for the target operating conditions.

\begin{figure}[h]
\centering
\begin{tikzpicture}[
    box/.style={draw, thick, rounded corners=3pt, minimum width=2.5cm,
                minimum height=1.2cm, align=center, font=\small},
    arr/.style={->, thick}
]
\usetikzlibrary{positioning}

% Nodes
\node[box, fill=gray!10]    (input)    {Input frame};
\node[box, fill=orange!15]  (hsv)      [above right=0.6cm and 1.5cm of input]  {\textbf{HSV}\\Colour masks};
\node[box, fill=cyan!15]    (hough)    [below right=0.6cm and 1.5cm of input]  {\textbf{Hough}\\Circle fitting};
\node[box, fill=green!15, minimum width=3.0cm]
                             (ensemble) [right=1.5cm of hsv.east |- input]      {\textbf{Ensemble merge}\\Union-Find + NMS};
\node[box, fill=gray!10]    (output)   [right=1.5cm of ensemble]               {\textbf{Output}\\[4pt]\textcolor{red!80}{$\bullet$}~0.97\quad\textcolor{blue!80}{$\bullet$}~0.95};

% Arrows
\draw[arr] (input.east) -- ++(0.5,0) |- (hsv.west);
\draw[arr] (input.east) -- ++(0.5,0) |- (hough.west);
\draw[arr] (hsv.east)   -| (ensemble.north);
\draw[arr] (hough.east) -| (ensemble.south);
\draw[arr] (ensemble.east) -- (output.west);

\end{tikzpicture}
\caption[Ensemble detection pipeline: HSV and Hough merged via Union-Find and NMS.]{Ensemble detection pipeline. HSV colour segmentation and Hough Circle
Transform independently process the same input frame. Their detections are merged
using a Union-Find structure and filtered by Non-Maximum Suppression (NMS).
Detections confirmed by both methods receive a confidence boost of 8\%, producing
the final output with associated confidence scores.}
\label{fig:ensemble}
\end{figure}


\hmsubsubsection{SVM Colour Verification}{OAH}{}
A Support Vector Machine (SVM) classifier was implemented as an additional colour
verification step following ensemble merging. The purpose of this secondary check
was to guard against cases where HSV thresholding misclassifies a ball due to
unusual lighting conditions or surface reflections.

The classifier is implemented in the \texttt{ColorHistogramClassifier} class and
operates on a normalised HSV histogram extracted from the bounding region of each
detected ball. The histogram consists of 100 features in total: 36 bins for the
hue channel over the range $[0, 180]$, 32 bins for the saturation channel, and
32 bins for the value channel. Pixels with saturation below 40 or value outside
the range $[30, 250]$ are excluded from the histogram to suppress background and
specular highlight pixels. Each channel histogram is independently normalised so
that its bins sum to 1.

The underlying model is a scikit-learn SVM pipeline stored as a 32~KB
\texttt{.pkl} file, requiring no TensorFlow dependency. Inference time on the
Raspberry Pi~5 is below 1~ms, making it suitable for real-time operation.

During system evaluation, the SVM was found to consistently misclassify dark
navy-blue balls as red with high confidence. This was attributed to the overlap
between dark blue and dark red in the HSV histogram space under low-saturation
conditions. As the HSV and Hough ensemble alone demonstrated sufficient colour
accuracy across all tested conditions, the SVM verification step was disabled.
The classifier infrastructure remains in the codebase for potential future
re-enablement should operating conditions change.

\hmsubsubsection{Object Tracking}{OAH}{}

To maintain stable object identities across consecutive frames, the system employs
a centroid-based multi-object tracker implemented in the BallTracker class. The
tracker assigns each detected ball a persistent integer identifier that remains
consistent as long as the object is visible, enabling downstream components to
refer to individual balls without ambiguity.

Each tracked object is associated with a Kalman filter operating on the state
vector $[x,\, y,\, v_x,\, v_y]^T$, using a constant-velocity motion model. The
transition and measurement matrices are defined as:
\begin{equation}
    F = \begin{pmatrix} 1 & 0 & 1 & 0 \\ 0 & 1 & 0 & 1 \\
    0 & 0 & 1 & 0 \\ 0 & 0 & 0 & 1 \end{pmatrix}, \qquad
    H = \begin{pmatrix} 1 & 0 & 0 & 0 \\ 0 & 1 & 0 & 0 \end{pmatrix}
\end{equation}
The process noise covariance is set to $Q = \mathrm{diag}(1.0,\, 1.0,\, 4.0,\,
4.0)$ and the measurement noise covariance to $R = \mathrm{diag}(25.0,\, 25.0)$.
These values reflect that the moment-based centroid measurement carries an
uncertainty of approximately $\pm 5$ pixels, while the velocity may change by up
to $\pm 2$ pixels per frame. The resulting Kalman gain is approximately 0.17 for a
stationary ball, meaning the corrected position is weighted 17\% towards the new
measurement and 83\% towards the predicted state.

At each frame, the tracker first advances all filters one timestep using the
prediction step. New detections are then matched to existing tracked objects by
constructing a cost matrix based on Euclidean distance between predicted positions
and measured centroids. Detections of different colours are assigned an infinite
cost and are never matched across colour classes. Matching is performed greedily in
order of ascending cost, and pairs whose distance exceeds the maximum threshold of
100 pixels are not matched. The time complexity of this procedure is $O(n^2)$,
which is sufficient for the expected maximum of six simultaneously tracked objects.

Matched filters are corrected using the new measurement. Unmatched existing tracks
advance via the predicted state without correction, a mode referred to as coasting.
If an object remains unmatched for more than 2 consecutive frames, it is
deregistered and its identifier is retired. Ball radius is smoothed using an
exponential moving average with a coefficient of $\alpha = 0.15$ to reduce
frame-to-frame jitter.
